\section{Failure Models in Distributed Systems}

Classical distributed systems theory adopts a taxonomy of failure modes: crash faults (a component halts), omission faults (a component fails to send or receive messages), timing faults (a component violates timing assumptions), and Byzantine faults (a component exhibits arbitrary behavior). The practical relevance of this taxonomy has been reevaluated in light of empirical studies of real-world failures.

Yuan et al.\ analyzed over 200 user-reported failures in distributed data-intensive systems and found that the majority of catastrophic failures originated from trivial mistakes in error-handling code rather than from the complex fault models studied in the theoretical literature. This finding motivates our focus on mechanisms that are robust to operator error and configuration drift, not only to well-modeled component crashes.
